\par\Needspace{6\baselineskip}\medskip
\begin{center}\hypertarget{native-174-article128}{}\pdfbookmark[2]{128.}{article-128}128.\end{center}
Wir bezeichnen die drei wie in den Artikeln 125 und 126 ermittelten Radiusvektoren mit $r,r',r''$.
Die heliozentrische Winkelbewegung in der Bahn vom zweiten zum \SourceLeaf{175}dritten Ort sei $2f$, die vom
ersten zum dritten $2f'$ und die vom ersten zum zweiten $2f''$, sodass $f'=f+f''$ gilt. Ferner gelte
\[
r'r''\sin2f=n,\qquad rr''\sin2f'=n',\qquad rr'\sin2f''=n'';
\]
und die Produkte der Konstanten $k$ (Artikel 2) mit den Zeitintervallen von der zweiten bis zur
dritten, von der ersten bis zur dritten und von der ersten bis zur zweiten Beobachtung bezeichnen
wir in dieser Reihenfolge mit $\theta,\theta',\theta''$. Wie im vorigen Artikel setzt man zur Bestimmung der Elemente
zwei Rechnungen an, die erste mit $r,r',f''$ und $\theta''$, die zweite mit $r',r'',f,\theta$. In beiden Rechnungen geht man
nur so weit, bis die Größe ermittelt ist, die das Verhältnis des elliptischen Sektors zum Dreieck
ausdrückt und oben (Artikel 91) mit $y$ oder $-Y$ bezeichnet wurde. Der Wert dieser Größe sei in
der ersten Rechnung $\eta''$, in der zweiten $\eta$. Nach Formel 18 in Artikel 95 erhält man so für den
Halbparameter $p$ zwei Werte
\[
\BareRoot{p}=\frac{\eta''n''}{\theta''}\qquad\text{und}\qquad
\BareRoot{p}=\frac{\eta n}{\theta}\mathpunct{\cdot}
\]
Nach Artikel 82 erhalten wir außerdem einen dritten Wert
\[
p=\frac{4rr'r''\sin f\sin f'\sin f''}{n-n'+n''}\mathpunct{\raisebox{.45ex}{,}}
\]
und alle drei Werte müssten offenbar übereinstimmen, wenn von Anfang an die wahren Werte für $x,y$
angesetzt worden wären. Es müsste daher gelten
\[
\frac{\theta''}{\theta}=\frac{\eta''n''}{\eta n}
\]
\[
\begin{aligned}
n-n'+n''&=\frac{4\theta\theta''rr'r''\sin f\sin f'\sin f''}{\eta\eta''nn''}\\[5pt]
&=\frac{n'\theta\theta''}{2\eta\eta''rr'r''\cos f\cos f'\cos f''}\mathpunct{\cdot}
\end{aligned}
\]
Sind diese Gleichungen nicht schon in der ersten Rechnung von selbst erfüllt, so kann man setzen
\[
\begin{aligned}
X&=\log\frac{\eta n\theta''}{\eta''n''\theta}\\[5pt]
Y&=n-n'+n''-\frac{n'\theta\theta''}{2\eta\eta''rr'r''\cos f\cos f'\cos f''}\mathpunct{\cdot}
\end{aligned}
\]

Dieses Verfahren ist ebenso allgemein anwendbar wie das zweite, das in Artikel 125 erläutert wurde.
Der große Vorteil dieses fünften Verfahrens besteht darin, dass die Rechnung bei den ersten
Hypothesen gleichsam auf halbem Weg zur Bestimmung der Elemente selbst stehen bleiben kann. Sobald
die Rechnung so weit gediehen ist, dass man absehen kann, dass die neue Hypothese von der Wahrheit
nicht merklich abweicht, genügt es, die Elemente selbst für diese Hypothese allein aus $r,r',f'',\theta''$, aus
$r',r'',f,\theta$ oder, vorzugsweise, aus $r,r'',f',\theta'$ zu bestimmen.
